\section{Introduction}\label{introduction}

Encrypted sparse linear algebra sits at an awkward boundary between cryptography and data layout. Homomorphic encryption hides values but makes data movement and permutation expensive. Sparse formats reduce arithmetic on zeros, yet they also turn a matrix's support into layout metadata and require a client to align the encrypted query with physical value lanes. CSSC addresses the static case by sorting rows, packing capacity-fitting rectangles, and retaining the metadata needed to recover logical outputs. Its own stated limitation is a static sparse pattern \cite{gao2026cssc}.

A mutable matrix cannot be handled by changing values alone. An insertion may consume padding, create an overflow component, change a row permutation, or force new query-reorganization metadata. A deletion may leave reusable physical capacity without changing the logical matrix dimension. Meanwhile, a query must not combine a new matrix component with an old column-index array or reconstruct shares using a plan from another version. These are system-level consistency requirements, not details of the underlying homomorphic primitive.

The naive policy republishes a fully recompressed CSSC layout after every batch. Alternative policies retain slack, patch local regions, append a delta, or periodically fold auxiliary state into the base. Their relative merit depends on the update/query ratio, freshness deadline, layout history, communication bandwidth, and the exact encrypted output path. Evaluating each window from a fresh initial layout, selecting with held-out information, or omitting returned ciphertexts and masking work can confound a comparison.

This paper therefore asks a narrower, auditable question: how should mutable CSSC state be published, queried, reconstructed, and compared without breaking version semantics or overstating evidence?

\subsection{Contributions}\label{contributions}

\begin{enumerate}
\def\labelenumi{\arabic{enumi}.}
\tightlist
\item
  \textbf{Version-bound publication semantics.} We define Publication Windows and a commit rule that binds logical state, CSSC components, global-column query metadata, the reconstruction plan, and prepared queries to one version.
\item
  \textbf{Explicit multi-component reconstruction.} We introduce a private RowMap-sensitive \texttt{OutputPlan} that distinguishes overlap summation, disjoint block concatenation, and implicit zeros.
\item
  \textbf{Scoped masking integration.} We derive overlap-only F1-M operands from the plan and bind every random mask to a persistent, reserve-before-sample five-field identity. The underlying canceling-mask idea is prior art; the contribution claimed here is its versioned CSSC integration.
\item
  \textbf{A designated strong fixed-segment reference path.} We define public fixed-width segment/page shapes, cloud-side leader reduction, private leader merging, and a deterministic whole-query execution bundle. We do not claim a new HE primitive or a formal security proof.
\item
  \textbf{Causal, fail-closed evaluation.} We freeze three strategy identities rather than train an online selector, advance independent states through identical deterministic streams, separate measured and projected quantities, and bind every reportable claim to independently replayed, digest-addressed evidence. A separate current-source study closes a 54-cell semantic-preservation, compile-call, and validation-lifecycle matrix without reopening the stopped strategy-cost campaign.
\end{enumerate}

\section{Background and Related Work}\label{background-and-related-work}

\subsection{Static encrypted sparse matrix--vector multiplication}\label{static-encrypted-sparse-matrixvector-multiplication}

CSSC is the direct substrate of this work and must be credited for row sorting, left-aligned sparse rectangles, \texttt{Value}, global \texttt{ColumnIndex}, \texttt{RowMap}, chunked query reorganization, and static aggregation \cite{gao2026cssc}. Our implementation is an independent pseudocode-derived reconstruction; we did not use or verify an author implementation. Static encrypted sparse linear algebra already spans Lodia's low-diagonal SpMV \cite{yu2025lodia}, CipherSkip's encrypted-value-and-index arbitrary-shape SpGEMM \cite{xiong2026cipherskip}, diagonal-packing reordering \cite{mutluergil2026diagonal}, ciphertext--ciphertext SpMSpM on CPU and GPU \cite{ferguson2025unstructured,dagata2026gpu}, and the encrypted-index Scatter--Gather--Apply design reported by the public SparseE program abstract \cite{wei2026sparsee}. Rhombus separately shows plaintext-matrix/encrypted-vector two-party MVM with additive output shares \cite{he2024rhombus}. These works rule out broad claims of first encrypted indices, first sparsity-aware FHE multiplication, first ciphertext--ciphertext sparse multiplication, or first use of random sharing to split an MVM output. Our claimed gap is the narrower mutable CSSC publication and reconstruction contract; it is not a new static packing or cryptographic primitive.

Our source audit of CSSC \cite{gao2026cssc} finds a non-power-of-two ambiguity in the printed aggregation pseudocode. We use a corrected \texttt{totalSum}-compatible stored-power/prefix graph that preserves the paper-derived abstract count \texttt{floor(log2\ w)\ +\ popcount(w)\ -\ 1}. We present this as a compatibility correction, not as a new aggregation algorithm or verified author-code behavior; the \texttt{totalSum} connection is grounded in the HElib SIMD reduction literature \cite{halevi2014algorithms}.

The closest methods are not interchangeable experimental baselines.

\textbf{Table 1. Closest method families and the boundary of this work.}

{\small\def\LTcaptype{none} % do not increment counter
\begin{longtable}[]{@{}
  >{\raggedright\arraybackslash}p{(\linewidth - 6\tabcolsep) * \real{0.2308}}
  >{\raggedright\arraybackslash}p{(\linewidth - 6\tabcolsep) * \real{0.2308}}
  >{\raggedleft\arraybackslash}p{(\linewidth - 6\tabcolsep) * \real{0.3077}}
  >{\raggedright\arraybackslash}p{(\linewidth - 6\tabcolsep) * \real{0.2308}}@{}}
\toprule\noalign{}
\begin{minipage}[b]{\linewidth}\raggedright
Work
\end{minipage} & \begin{minipage}[b]{\linewidth}\raggedright
Principal representation or setting
\end{minipage} & \begin{minipage}[b]{\linewidth}\raggedleft
Mutable support?
\end{minipage} & \begin{minipage}[b]{\linewidth}\raggedright
Relation to this paper
\end{minipage} \\
\midrule\noalign{}
\endhead
\bottomrule\noalign{}
\endlastfoot
CSSC \cite{gao2026cssc} & Sparse-coordinate compression with \texttt{ColumnIndex} / \texttt{RowMap} & No & Direct static substrate; its published layout and reorganization are not our contribution. \\
Lodia \cite{yu2025lodia} & Low-diagonal decomposition for batched FHE SpMV & Not described & Establishes earlier sparsity-aware encrypted SpMV; different decomposition and leakage/cost interface. \\
CipherSkip \cite{xiong2026cipherskip} & Encrypted values and indices for arbitrary-shape SpGEMM and chained products & No matrix-state publication protocol & Its server-side alignment concerns encrypted intermediate products, not insert/delete/update of a published sparse layout. \\
SparseE \cite{wei2026sparsee} & Encrypted-index Scatter--Gather--Apply with a permutation/expansion accelerator & Not established by the public descriptions & Establishes an encrypted-index hardware co-design boundary; as of 2026-08-30, an official program abstract and a BUAA institutional summary were public, but no public full text, DOI, or public software repository was located. \\
Diagonal Packing \cite{mutluergil2026diagonal} & Row/column reordering to reduce occupied cyclic diagonals & Not described & HE-aware static compiler optimization; no mutable CSSC publication or reconstruction is described. \\
Ferguson et al. \cite{ferguson2025unstructured} & CKKS ciphertext--ciphertext SpMSpM with public sparse metadata & Not described & Demonstrates prior FHE sparsity exploitation under a PPML workload; different operation, scheme, and small square matrices. \\
D'Agata et al. \cite{dagata2026gpu} & GPU/FIDESlib CKKS ciphertext--ciphertext SpMSpM with public sparse metadata & Not described & Extends the same static design space to GPU execution; a dynamic SpMV publication protocol is not described. \\
Rhombus \cite{he2024rhombus} & Plaintext-matrix/encrypted-vector MVM with additive output shares & Not described & A different two-party PPML protocol; the random-share primitive is prior art, while versioned overlap binding is our narrower integration claim. \\
Damie et al. \cite{damie2025secure} & Secret-shared sparse matrix multiplication in MPC & Not described & Shows secure sparse arithmetic outside FHE; incomparable without a common threat and cost model. \\
This work & Version-bound CSSC base plus designated maintenance components; comparative evaluation stopped, separate validation scaling completed & Yes & Studies publication, query recompilation, private reconstruction, and causal evidence contracts; claims no static-format, primitive, or comparative-performance novelty. \\
\end{longtable}
}

\subsection{Mutable sparse layouts}\label{mutable-sparse-layouts}

Reserved slack, padding reuse, COO/HYB overflow, row-local growth, base/delta organization, periodic compaction, and format selection all have substantial plaintext sparse-matrix and storage-system precedent, including SELL-C-sigma \cite{kreutzer2014sellcsigma}, GPU HYB/COO practice \cite{bell2008spmvcuda}, LSM-trees \cite{oneil1996lsm}, Dynamic-CSR \cite{king2016dynamiccsr}, and dynamic format selection \cite{stylianou2022morpheus}. They are comparison families in this paper, not independent novelty claims. Our narrower concern is their interaction with encrypted query preparation, version consistency, private reconstruction, and complete cost accounting.

\subsection{Dynamic encrypted data systems and masking}\label{dynamic-encrypted-data-systems-and-masking}

Dynamic SSE, encrypted databases, ShieldDB, and oblivious transaction systems already study updates, client state, epochs, padding, and leakage \cite{cash2014dynamic,vo2021shielddb,crooks2018obladi}. They do not instantiate the same ciphertext--ciphertext CSSC query, but they prevent broad claims that batching or versioned updates are new. d-DSE directly shows that update-volume leakage remains a separate problem and that padding can impose large storage and communication costs \cite{liu2024ddse}. CKKS-Auth Tree uses versioned root commitments and timestamps to detect stale or replayed verification objects after encrypted-metadata updates \cite{chen2026ckksauthtree}. Neither work implements mutable SpMV, but together they rule out novelty claims for volume hiding, versioned commitments, freshness checks, or replay rejection in isolation. Fixed segments, dummy work, and visible schedules are therefore described through an explicit leakage surface, not as a generic update-leakage solution.

Canceling one-time masks also predate this work \cite{bonawitz2017secureagg}. We therefore state the exact leakage and recipient roles and avoid the terms \texttt{simulation-secure} or \texttt{prevents\ leakage} without a separate proof. Authorized homomorphic database updates also predate this work and delimit any claim about encrypted mutability \cite{parbat2023authorized}.

\section{System and Threat Model}\label{system-and-threat-model}

The system has Client A, which owns the mutable matrix, updates, CSSC metadata, component \texttt{ColumnIndex} arrays and \texttt{RowMap}s, the complete \texttt{OutputPlan}, and the F1-M reservation/sampling ledger; Client B, which owns the query and secret key and receives the result; and a semi-honest Cloud evaluator. Client A and Client B are both authorized to read the complete \texttt{RowMap}-sensitive \texttt{OutputPlan}; \texttt{private} means private from the Cloud, not private between the two clients. Client A sends the versioned column metadata and complete plan to Client B, samples and encrypts F1-M operands under Client B's public key, and never sends the individual mask plaintexts to Client B. At most one party is corrupted and the Cloud does not collude with either client. The model excludes malicious behavior, adaptive corruption, availability, side-channel, traffic-analysis, and collusion claims.

Client B is authorized to learn the versioned global column indices, component \texttt{RowMap}s, complete reconstruction plan, and final result. The Cloud is authorized to learn public parameters, ciphertext shapes and counts, public page/segment shapes, the operation schedule, opaque identifiers, query/version identifiers, and plan digests. The interface classifies matrix/query plaintexts, global column metadata, RowMaps, full plans, mask plaintexts, and unblinded component outputs as forbidden-to-Cloud fields, and the typed serializers do not emit them. This is an access-control and serialization requirement, not a proof that all leakage is prevented or that the implementation realizes a simulation-based notion.

\section{Design}\label{design}

Figure 1 shows the protocol boundary. Client A owns matrix publication, the complete plan and private metadata, constructs the version commit, and prepares the encrypted F1-M operands. Client B owns the query and key set, uses the version-matched column metadata to gather and encrypt query values, and uses its copy of the complete plan to recover logical coordinates. The Cloud receives only the typed public program, encrypted operands, and permitted identifiers. The diagram is an interface summary, not a claim that the listed metadata is cryptographically hidden from every side channel.

\begin{figure}
\centering
\includegraphics[width=\linewidth,height=\textheight,keepaspectratio]{assets/route-c-protocol-flow.png}
\caption{Version-bound mutable CSSC publication, execution, and private reconstruction.}
\end{figure}

\subsection{Publication Windows}\label{publication-windows}

Updates accumulate until a query arrival, freshness deadline, microbatch bound, or explicit publication event closes the current window. The transition first applies every net update to a candidate logical state, validates matrix and value bounds, constructs or patches physical components, decodes them back to the same logical state, and only then publishes the next version. A failed transition cannot partially advance the visible state.

For the publication experiment, accepted raw events form atomic scheduling groups. All visible SET transitions are applied first, followed by a payload-free logical clock TICK and then the group's exactly scheduled queries. The TICK is emitted even for clipped no-ops, so freshness advances without fabricating an update. Before a group at logical time \texttt{t} is processed, any pending half-open window whose deadline is at or before \texttt{t} is closed; the deadline group therefore belongs to the next window. No close may split the group's SET→TICK→query sequence. One deterministic bound-one ternary query vector is frozen per paired trace unit and reused across candidates and cells. Its seed, generation method, and values are public, so it is a known-answer correctness/cost control rather than evidence for a production-query distribution, cryptographic randomness, or query-plaintext confidentiality.

\subsection{Versioned query reorganization}\label{versioned-query-reorganization}

Each CSSC value lane is paired with a parallel original matrix-column identifier or a padding sentinel. Client A sends versioned plaintext column metadata to Client B. Client B gathers one aligned vector per encrypted value chunk and encrypts it. The matrix column domain is independent of the ciphertext-slot domain; reducing a global column identifier modulo the slot count changes the function and is forbidden.

\subsection{OutputPlan and F1-M}\label{outputplan-and-f1-m}

The plan maps \texttt{(component,\ output\ block,\ physical\ lane)} to a logical output coordinate. Client B initializes the public-length result to zero, reorders each decrypted share, sums only equal logical coordinates, and concatenates disjoint horizontal blocks. A coordinate with no physical contributor is an implicit zero.

For contributor multiplicity greater than one, Client A samples a modular zero-sum mask tuple. The binding \texttt{(query,\ version,\ plan\ digest,\ component,\ output\ block)} is reserved in the durable SQLite ledger before sampling. Duplicate use is rejected within that ledger model. Repository rollback, database cloning or compromise, and cross-device coordination are outside the present evidence. A disjoint strong-path return uses an encrypted-zero dummy to retain a uniform visible operation position; it is not a random mask and is accounted separately.

\subsection{Strong fixed-segment delta}\label{strong-fixed-segment-delta}

The strong path retains a real CSSC base and places overflow into fixed-width power-of-two segments, currently frozen to \texttt{c=128}. A segment contains entries from one logical row. The public typed plan uses uniform shapes and opaque ordinal IDs, while the Cloud-visible leakage also includes published counts, schedule and timing, query/version identifiers, and binding digests. We do not claim segment unlinkability. The Cloud multiplies aligned value/query ciphertexts, applies the fixed reduction schedule, selects segment leaders, adds one F1-M operand per return, and returns page shares. Client B uses the private plan to merge leaders that map to the same logical row.

The current policy does not authorize cloud-side merging of segments that share a logical row. Same-row-equivalence fields are carried only in Client B's typed plan; this is an interface access-control requirement, not a proof that segment relationships cannot be inferred. The publication identity for folding/compaction is frozen separately; changing it creates a different candidate.

The segment width 128 is fixed before real-stream evaluation rather than tuned: it yields a seven-stage power-of-two reduction, places 32 complete segments in the 4,096-lane effective domain, and is the exact 127-active-plus-one-padding boundary covered by the pinned witness. We make no optimality or cross-segment-width claim.

\subsection{Definition-level functional propositions}\label{definition-level-functional-propositions}

The following statements are conditional functional propositions about the frozen typed interfaces and algorithms. They are not empirical conclusions and not a malicious- or simulation-security theorem. The exact-S1 source-conformance record maps every premise below to the corresponding validator, state transition, reconstruction routine, or ledger transition; the registered tests exercise legal cases and substitutions but do not replace the proof arguments.

\subsubsection{P1: binding soundness}\label{p1-binding-soundness}

For publication version \(v\), let

\[
\mathcal{P}^{(v)}=
\left(
s^{(v)},
\{C_k^{(v)}\}_{k=1}^{K_v},
\{CI_k^{(v)}\}_{k=1}^{K_v},
\{RM_k^{(v)}\}_{k=1}^{K_v},
O^{(v)},Q^{(v)},B^{(v)}
\right),
\]

where \(s^{(v)}\) is the canonical logical-state identity binding the logical matrix \(A^{(v)}\), \(C_k\) are ordered physical components, \(CI_k\) and \(RM_k\) are their global-column and row mappings, \(O\) is the private OutputPlan, \(Q\) is the prepared-query binding, and \(B\) is the immutable typed identity and digest bundle. Let \(B^{(v,q)}\) be the authoritative binding for query \(q\). The registered acceptance predicate has the shape

\[
\operatorname{Accept}(X;B^{(v,q)})=
\bigwedge_{f\in\mathcal{F}}[X_f=B_f^{(v,q)}]
\;\land\!
\bigwedge_{o\in\mathcal{O}}
[\operatorname{SHA256}(\operatorname{bytes}_X(o))=B_o^{(v,q)}],
\]

where \(\mathcal{F}\) enumerates every typed version, query, component, mapping, plan, parameter, and execution identity, and \(\mathcal{O}\) enumerates every retained canonical byte object. Hence acceptance implies equality for every enumerated field and rehashed object. The capability and formal-artifact paths are downstream of this conjunction, so a failed conjunct cannot mint them.

The proof is a case analysis over the closed predicate: each field class occurs as a necessary conjunct, retained bytes are rehashed inside the private admission process, and no earlier branch reaches the issuer. This establishes a functional binding property for ordinary inputs and malformed/stale substitutions. It says nothing about deliberate SHA-256 collisions, compromised hosts, or side channels.

\subsubsection{P2: multi-component reconstruction}\label{p2-multi-component-reconstruction}

Let \(A_k^{(v)}\) be the physical-output-row by global-logical-column matrix represented by component \(k\), and let \(E_k\) embed its declared physical output rows into the logical matrix domain. The admitted decomposition is complete when

\[
A^{(v)}=\sum_{k=1}^{K_v}E_k\!\left(A_k^{(v)}\right)\pmod t.
\]

Assume each component execution decrypts to \(A_k^{(v)}x\) in its declared lanes, every lane has exactly its declared logical coordinate, the OutputPlan is total, and P3 masks cancel. For logical row \(r\), the reconstruction operator sums all contributors mapped to \(r\), places disjoint blocks in logical order, and uses zero when the contributor set is empty. Therefore

\[
R_{O^{(v)}}\!\left(\{A_k^{(v)}x\}_{k=1}^{K_v}\right)_r
=\sum_{k=1}^{K_v}\left(E_k(A_k^{(v)})x\right)_r.
\]

By completeness of the admitted decomposition,

\[
\sum_{k=1}^{K_v}\left(E_k(A_k^{(v)})x\right)_r
=\left(A^{(v)}x\right)_r\pmod t.
\]

The equality is coordinatewise: overlap is ordinary addition in \(\mathbb{Z}_t\), disjoint blocks have disjoint images under the embeddings, and an unmaterialized row contributes the additive identity. It holds for any admitted number of components under the stated completeness and totality premises; the finite tests only check the implementation of those premises.

\subsubsection{P3: mask cancellation and ledger-scoped no-reuse}\label{p3-mask-cancellation-and-ledger-scoped-no-reuse}

For an overlap group \(G_r=(k_1,\ldots,k_g)\) in canonical contributor order with \(g\ge 2\), Client A samples

\[
m_{r,k_1},\ldots,m_{r,k_{g-1}}
\overset{\mathrm{u.a.r.}}{\leftarrow}\mathbb{Z}_t,
\qquad
m_{r,k_g}=-\sum_{i=1}^{g-1}m_{r,k_i}\pmod t.
\]

Thus \(\sum_{k\in G_r}m_{r,k}=0\pmod t\), and P2 reconstruction removes the random operands without changing the logical output coordinate. Groups of size zero or one receive no random tuple; an encrypted-zero dummy is a separate accounted object.

For no-reuse, the durable state machine is

\[
\textsf{unseen}\rightarrow\textsf{reserved}
\rightarrow\textsf{prepared}\rightarrow\textsf{consumed}.
\]

One transaction reserves every five-field key \texttt{(query,\ version,\ plan\ digest,\ component,\ output\ block)} before the first random sample. A prepared batch additionally binds the private plan, execution binding, modulus, operand commitments, and a unique token; verification consumes that token atomically. Uniqueness constraints plus exact terminal closure reject a duplicate key or token, commitment drift, orphan record, or second consumption. The evaluation-lane digest includes the unit-attempt ordinal, so the single provider-replacement allowance cannot reuse the same reservation identity. This is an invariant of one uncompromised durable SQLite ledger, not a rollback, cloning, compromise, or cross-device guarantee.

\subsubsection{P4: fixed-segment reconstruction}\label{p4-fixed-segment-reconstruction}

For a logical row \(r\), partition its auxiliary post-multiplication lane values \(a_{r,c}x_c\) in canonical entry order into \(J_r\) segments \(S_{r,1},\ldots,S_{r,J_r}\) of \(c=128\) lanes, so every populated \(S_{r,j}[\ell]\) is one such \(a_{r,c}x_c\) value, and pad only the final unused suffix with zeros. The seven-stage public reduction places

\[
L_{r,j}=\sum_{\ell=0}^{127}S_{r,j}[\ell]
\]

in the predetermined leader lane. The private OutputPlan maps every leader back to \(r\), so Client B obtains

\[
\sum_{j=1}^{J_r}L_{r,j}
=\sum_{j=1}^{J_r}\sum_{\ell=0}^{127}S_{r,j}[\ell],
\]

which is exactly the auxiliary contribution to row \(r\); adding the CSSC-base contribution and applying P2 gives the direct logical product. The proof is an induction on \(J_r\): one segment is the fixed reduction identity, and appending a row-owned segment adds exactly its leader sum without changing prior leaders. The 127/128/129, tombstone, padding, disjoint, and overlap cases are boundary tests of this construction, not its general proof and not evidence that 128 is an optimal width.

\section{Implementation and Evidence}\label{implementation-and-evidence}

At historical baseline commit \texttt{blinded-baseline}, the Python implementation supplied typed persistent states, a common query compiler, exact operation graphs, private operand/route bindings, a plaintext oracle, SQLite no-reuse commitments, canonical artifacts, and a separate deterministic replay validator. The encrypted path uses BFV-family batching, whose original scheme lineage is Fan--Vercauteren \cite{fan2012somewhat}, through OpenFHE 1.5.1 pinned by source commit \cite{openfhe151}.

At that historical baseline, R0 passed 750 tests and the Phase 2 whole-query witness passed in run \texttt{blinded-baseline-run}. The witness exercises a 4096-by-8193 fixture, global-column anti-aliasing, the 127-of-128 segment boundary, the unused second BFV batching row, explicit no-relinearize/relinearize transitions, random and dummy F1-M operands, and equality with both typed and direct plaintext oracles. Its GitHub Actions artifact-wrapper SHA-256 digest is recorded in the version-bound claim-ledger supplement.

A separate preregistered current-source replication retained that fixture,
expected output, OpenFHE pin, and witness specification while binding them to
reviewed source \texttt{blinded-E4-source}. Its sole run and 19-file artifact
were independently rehashed against a detached source checkout. All 35 records
in the eight-case deterministic contract passed, and the fixed whole-query
witness decrypted to \texttt{[(0, 128), (4095, 5)]}, matching both plaintext
oracles. Section~\ref{current-source-deterministic-conformance-outcome} states
the result and its exclusions.

A second current-source lineage froze a two-operation producer/replay interface
for validation-lifecycle scaling. One tagged attempt created three private producer
handoffs, three redacted independent-replay packages, and one aggregate. The
aggregate rehashed every provider object, closed all 54 cells, and was rebuilt
byte-for-byte in a separate post-terminal audit. Section~\ref{current-source-validation-lifecycle-scaling}
states the method; Section~\ref{current-source-validation-scaling-outcome}
states the bounded result.

This evidence is deliberately narrow. It does not establish the stopped
campaign's complete candidate costs, candidate registration, mixed-circuit
parameter safety, security, native performance, or end-to-end deployment.

\textbf{Table 2. Definition-level obligations and their Route C disposition.}

{\small\def\LTcaptype{none} % do not increment counter
\begin{longtable}[]{@{}
  >{\raggedright\arraybackslash}p{(\linewidth - 4\tabcolsep) * \real{0.3333}}
  >{\raggedright\arraybackslash}p{(\linewidth - 4\tabcolsep) * \real{0.3333}}
  >{\raggedright\arraybackslash}p{(\linewidth - 4\tabcolsep) * \real{0.3333}}@{}}
\toprule\noalign{}
\begin{minipage}[b]{\linewidth}\raggedright
Obligation
\end{minipage} & \begin{minipage}[b]{\linewidth}\raggedright
Frozen implementation and review boundary
\end{minipage} & \begin{minipage}[b]{\linewidth}\raggedright
Claim permitted in this manuscript
\end{minipage} \\
\midrule\noalign{}
\endhead
\bottomrule\noalign{}
\endlastfoot
P1: exact binding & Typed version/query/plan/payload identities, fresh byte rehashing, independent replay, and source-conformance mapping & Conditional binding proposition and fail-closed interface design; no end-to-end admission claim \\
P2: multi-component reconstruction & Private \texttt{Output}\allowbreak\texttt{Plan}, direct plaintext oracle, overlap/\allowbreak concatenation/\allowbreak implicit-zero checks & Definition-level reconstruction proof under the admitted-plan premises; no formal-run result \\
P3: F1-M cancellation and no-reuse & Modular cancellation proof plus durable reserve-before-sample SQLite state machine & Cancellation and ledger-local invariant only; no simulation-security, rollback, cloning, or cross-device claim \\
P4: fixed 128-lane segment path & Seven-stage reduction proof and 127/128/129 boundary tests & Correctness proposition for the frozen segment construction; no optimal-width or performance claim \\
Evidence release & S1/S2 Behavior Sets, independent guards, terminal controller, exact object-count contracts, and a distinct validation-scaling source tag & Engineering/source conformance, the Route C stop, and bounded current-source semantic/compile/lifecycle observations; no strategy-cost evidence \\
\end{longtable}
}

The stopped Route A lineage used a permanently non-admissible six-job
qualification before a conditional, strictly serial formal campaign. Exact S1
CI/PRE-S1 and the data-only S2 compatibility gate passed, but the sole
qualification missed its frozen deadline; no acquisition, formal unit,
terminal admission, aggregate, or analysis object was created. A later
follow-up preserved the matrix and thresholds and passed five authority-false
controls, but its one provider attempt stopped before checkout or seed
admission and produced zero artifacts. Both frozen one-shot rules therefore
closed without strategy-cost evidence or permission to rerun.

Generative-AI systems assisted literature discovery, code and test generation, adversarial review, and drafting during development. Their outputs are neither source authority nor experimental evidence. Human authors must independently verify cited primary sources, inspect the final code, reproduce the admitted artifacts, and accept responsibility for every claim and released result.

\section{Evaluation Methodology}\label{evaluation-methodology}

Sections 6.1--6.6 preserve the frozen counterfactual execution and reporting protocol. Because neither qualification path produced a GO, none of the acquisitions, formal cells, native cases, measurements, or reports specified below occurred. The protocol was frozen before any formal artifact could be inspected. It specified no trained selector, winner chosen on a tuning prefix, or held-out oracle. Every strategy would have owned an independent persistent state and consumed the same ordered event groups, publication boundaries, queries, initial logical matrix, and public parameters. The intended intervention was therefore the maintenance mechanism, not the workload history. Section~6.7 is a distinct functional replication. Section~6.8 is a separate current-source validation-scaling study. Neither rescues or estimates the stopped strategy-cost campaign.

\subsection{Fixed strategies}\label{fixed-strategies}

The three strategy identities were:

\begin{enumerate}
\def\labelenumi{\arabic{enumi}.}
\tightlist
\item
  \texttt{periodic-repack/windows=1}, which would have reconstructed and republished a complete static CSSC layout at every Publication Window;
\item
  \texttt{padding-reuse}, which would have reused the lowest-ordinal tombstone, then natural padding, and otherwise rebuilds the affected fixed horizontal row partition; and
\item
  \texttt{packed-coo-cloud-segmented-delta/segment-width=128}, which would have kept a CSSC base and appended overflow into cloud-executable, row-owned 128-lane segments. It was not to fold online in the registered experiment.
\end{enumerate}

These candidates were mechanisms specified for comparison, not actions of an online policy. A slower strategy, a crossing cost curve, or the absence of a global winner would have been a reportable outcome rather than a failed experiment.

\subsection{Synthetic matrix}\label{synthetic-matrix}

The synthetic suite was specified to use a mixed insert/delete/modify workload at two frozen scales: S would have had 256 rows, 8,193 columns, and 512 accepted updates; M would have had 1,024 rows, 8,193 columns, and 2,048 accepted updates. The formal seeds were \texttt{20260822}, \texttt{20260823}, and \texttt{20260824}, defining six scale--seed shards. Every shard would have evaluated all three strategies at \(\rho\in\{0.01,0.1,1,10\}\).

For a reduced fraction \(\rho=p/q\) and zero-based accepted-event ordinal \(a\), the number of queries inserted after event group \(a\) is

\[
Q_a=
\left\lfloor\frac{(a+1)p}{q}\right\rfloor-
\left\lfloor\frac{ap}{q}\right\rfloor,
\]

so the first \(N\) complete groups contain exactly

\[
\sum_{a=0}^{N-1}Q_a=\left\lfloor\frac{Np}{q}\right\rfloor
\]

queries. The \(0.01\), \(0.1\), and \(1\) cells were specified for direct execution. The \(\rho=10\) cell would not have been timed; it was registered as an exact transformation of the same strategy and shard's \(\rho=1\) query-linear count and byte fields. Any event, window, state, or non-whitelisted-field mismatch would have rejected the entire shard.

\subsection{Ordered events from one real source}\label{ordered-events-from-one-real-source}

The sole external object would have been the SNAP Stack Overflow answer-to-question stream \texttt{sx-stackoverflow-a2q.txt.gz} \cite{snap_stackoverflow_temporal_network,paranjape_motifs_in_temporal_networks}. The acquisition job was specified to record the final URL, response headers, exact compressed length, and SHA-256 digest. An independent guard would have downloaded the object again and required identical response-body bytes; an admitted artifact would have omitted the raw compressed object.

The first 1,000,000 eligible records were specified to freeze the row/column mapping. Source-ID hashing would have defined two deterministic partitions, each with 1,024 rows and 8,193 columns. After the mapping prefix, each partition would have consumed 4,096 accepted records under two semantics. T1 was defined as cumulative occurrence,

\[
A_{uv}(t)=\min\{7,N_{uv}(t)\},
\]

whereas T2 would have retained the most recent \(K=1024\) accepted events and performed expiry before admission in one indivisible atomic group. The synthetic logical clock would have advanced by one second per 128 accepted records. Both partitions and both semantics were specified to run all three strategies at \(\rho\in\{0.1,1\}\), yielding four ordered-event shards. If admitted, those shards could only have supported conclusions about deterministic interactions within one source, not multi-source or historical wall-clock generalization.

\subsection{Native OpenFHE cases}\label{native-openfhe-cases}

The native matrix was specified to contain three strategies at the S and M scales, hence six cases. Every case would have used seed \texttt{20260822}, \(\rho=1\), the terminal accepted-event prefix (512 for S and 2,048 for M), and the query after the last complete event group. Version, component inventory, query vector, \texttt{OutputPlan}, typed execution plan, and canonical input bytes were frozen together.

Each case was specified to perform one discarded warm-up, three fresh-key producer evaluations, and three exact package replays:

\[
6\times(1+3+3)=42
\]

native evaluations in total. The three producer repetitions specified for recording were technical repetitions, not independent population samples. An admitted execution would have reported their raw values, median, and range, separately for production and replay. Replay was specified to deserialize the retained context, secret/public keys, evaluation-key frame, and input ciphertexts. Its lifecycle receipt would have had to show zero context generation, zero key generation, zero evaluation-key generation, and zero encryption, while its cloud-program operation inventory would have had to equal that of the corresponding producer package.

\subsection{Measurements and accounting}\label{measurements-and-accounting}

For strategy \(k\), scale or source unit \(s\), and query/update ratio \(\rho\), the result was specified as a typed vector rather than a synthetic score:

\[
\mathbf{c}_{k,s,\rho}=
\left(
T_{\mathrm{state}},
T_{\mathrm{assembly}},
N_{\mathrm{op}},
B_{\mathrm{meta}},
\overline{B}_{\mathrm{crypto}},
\mathrm{RSS}_{\max},
\mathrm{scratch}_{\max}
\right).
\]

Synthetic and ordered-event cells were specified to measure state-transition, result-assembly, and independent-replay time, peak RSS, and controlled scratch; they would have exactly counted events, windows, queries, typed operations, and object multiplicities. Their cryptographic-object bytes were specified as conservative type-derived bounds,

\[
\overline{B}_{\mathrm{crypto}}=\sum_j m_jU_j,
\]

where \(m_j\) is the exact multiplicity and \(U_j\) is the registration-time, pre-execution Stage-1 maximum serialized size frozen for type \(j\). Canonical metadata bytes would have been measured separately. Native cases instead were specified to report observed serialized-object bytes, typed-operation inventories, process time, RSS, and scratch. Simulator counts were never to be converted into OpenFHE latency, and evidence-replay overhead was never to be presented as strategy execution cost.

For the registered \(\rho=10\) projection, only query-linear fields change:

\[
n^{(10)}_{q,p}=10n^{(1)}_{q,p},\qquad
n^{(10)}_{u,p}=n^{(1)}_{u,p}.
\]

Wall time, RSS, scratch, and native latency would have been unavailable for this projected cell. Protocol-byte transfer time at bandwidth \(b\) Mbps was specified only as the transparent conversion

\[
T_{\mathrm{net}}(B,b)=\frac{8B}{b\times10^6};
\]

HTTP/TLS, artifact wrappers, filesystems, and private replay transport were to remain separate evidence-pipeline costs.

The paired strategy contrast for unit \(u\) is

\[
\tau_{a,b}(u,\rho)=C_a(E_u,\rho)-C_b(E_u,\rho),
\]

where \(C\) is one preregistered field or the full typed cost vector. With only two scales, two deterministic source partitions, and three technical native repetitions, an admitted execution would have reported all raw points and mechanism-level decompositions. It would not have fitted scaling exponents, attached population \(p\) values or confidence intervals, or claimed a global winner or Pareto frontier.

\subsection{Evidence admission and stopping}\label{evidence-admission-and-stopping}

Before formal execution, one six-job qualification was required to complete within the frozen 45-minute computational path and 55-minute total path, satisfy the necessary \(6C_q\leq9000\,\mathrm{s}\) planning screen, and survive a fresh external-controller reread. Here \(C_q\) is the registered per-case native-runtime estimate: the sum of provider-API \texttt{startedAt}--\texttt{completedAt} durations for the qualification's native producer job, native replay/guard job, and entire combined-guard job. Qualification artifacts were permanently non-evidence and could not enter the formal result set.

The formal campaign was specified as serial and would have admitted exactly

\[
1\ \text{acquisition}
+6\ \text{synthetic}
+4\ \text{ordered-event}
+6\ \text{OpenFHE}
=17
\]

pre-aggregate artifacts. Before launching each unit, the controller would have reserved its full remaining worst-case budget. The preregistered 12-hour figure was an acceptance threshold, not a guarantee that provider-side cancellation cannot overshoot. Across the 17 enumerated units, at most one whole-unit replacement would have been allowed, and only for a provider failure rather than a data-dependent or scientific outcome. Terminal admission would have rejected missing, extra, duplicated, wrong-attempt, or wrong-kind objects before aggregation and compatible detached analysis.

\subsection{Current-source deterministic conformance replication}\label{current-source-deterministic-conformance-replication}

After both performance lineages were terminal, a separate preregistration froze
a functional question that neither estimates nor rescues their performance
estimand. One immutable reviewed source, one default dispatch, CPython 3.12, a
120-minute job timeout, two build threads, and OpenFHE 1.5.1 at source commit
\texttt{1306d14f8c26} were fixed before execution.

The deterministic corpus froze seed \texttt{20260822}, eight input cases, and
35 contract records. The encrypted witness froze one 4096-by-8193
CSSC-base-plus-strong-delta fixture, segment width 128, active payload 127,
padding offset 127, global column 8192, and expected sparse output
\texttt{[(0, 128), (4095, 5)]}. PASS required exact source/run/attempt identity,
success for every provider step, one success-named artifact, exactly 19 regular
files, raw-ZIP/provider-digest equality, strict internal checksums, literal
RFC~6901 checks, and detached-source provenance rehashing. No replacement was
needed or dispatched.

\subsection{Current-source validation-lifecycle scaling}\label{current-source-validation-lifecycle-scaling}

After both performance lineages were closed, a separate preregistration froze
a descriptive current-source question: how producer and independent-replay
validation-bearing lifecycle duration changes with exact query count while
producer/replay semantic projections are preserved and the same preregistered
compile-call gate \(Q\leq C\leq2Q\) is applied at each \(Q\). The S-scale matrix crossed
the three fixed strategy identities with \(\rho\in\{1/100,1/10,1\}\), hence
\(Q\in\{5,51,512\}\), three precommitted SHA-derived seeds, and two roles, for
54 cells. Cyclic rotations placed each strategy block and each query count in
each corresponding order position once across the three seeds.

The producer and replay were distinct deep operations in distinct hosted-runner
jobs. Wall and process clocks surrounded only the operation call; setup, trace
generation, provider transfer, and queue time were excluded. A shared wrapper
counted every \texttt{compile\_query} call in both live bindings, with the
preregistered gate \(Q\leq C\leq2Q\). Replay additionally required byte-exact
equality of the frozen semantic projection. One attempt had to create exactly
three private producer artifacts, three redacted replay artifacts, and one
aggregate; any failed job, extra attempt, incomplete cell, or mismatch made all
timing unreportable.

For each strategy and role, the analysis reports all three observations and
their integer-nanosecond median, minimum, and maximum at each \(Q\). It fits
\(T(Q)=\alpha+\beta Q\) by unweighted OLS to the three median wall times using
exact rational arithmetic. The fit has no acceptance threshold; \(\alpha\) is
an extrapolated descriptive intercept, not an observed fixed cost. Three
query counts do not support an asymptotic, population, or ranking claim.

\section{Evidence Outcome}\label{evidence-outcome}

\subsection{Evidence status}\label{evidence-status}

Exact-head CI, PRE-S1, descriptive registration, and S1--S2 compatibility gates closed before execution. Qualification run \texttt{blinded-qualification} was then dispatched once from exact S2 while importing behavior from exact S1. At the frozen 45-minute computational threshold, the independent simulator replay was still running and the combined guard had not started. The external controller requested cancellation of that exact run. The terminal workflow conclusion is \texttt{cancelled}; the only retained object was the one-day q1 pre-replay handoff, which is permanently non-admissible.

No primary q5 guarded bundle, q6 post-run record, live dispatch capability, acquisition artifact, formal shard, terminal admission record, Route A aggregate, or compatible S3 analysis exists. The project therefore has zero reportable Route A strategy-cost, ordered-event, or native-performance results. CI logs, qualification handoffs, partially completed jobs, and local copies are not alternative result sources.

Figure 2 records the evidence boundary rather than a performance curve. Green and blue boxes are engineering or workflow observations at their exact source identities; none is promoted into the absent formal result set.

\begin{figure}
\centering
\includegraphics[width=\linewidth,height=\textheight,keepaspectratio]{assets/route-c-evidence-boundary.png}
\caption{Primary exact-source freeze, one-shot qualification stop, and the resulting Route C evidence boundary; the later follow-up is recorded separately.}
\end{figure}

\textbf{Table 3. What each evidence layer proves and does not prove.}

{\def\LTcaptype{none} % do not increment counter
\begin{longtable}[]{@{}
  >{\raggedright\arraybackslash}p{(\linewidth - 6\tabcolsep) * \real{0.2500}}
  >{\raggedright\arraybackslash}p{(\linewidth - 6\tabcolsep) * \real{0.2500}}
  >{\raggedright\arraybackslash}p{(\linewidth - 6\tabcolsep) * \real{0.2500}}
  >{\raggedright\arraybackslash}p{(\linewidth - 6\tabcolsep) * \real{0.2500}}@{}}
\toprule\noalign{}
\begin{minipage}[b]{\linewidth}\raggedright
Layer
\end{minipage} & \begin{minipage}[b]{\linewidth}\raggedright
Exact disposition
\end{minipage} & \begin{minipage}[b]{\linewidth}\raggedright
Supports
\end{minipage} & \begin{minipage}[b]{\linewidth}\raggedright
Does not support
\end{minipage} \\
\midrule\noalign{}
\endhead
\bottomrule\noalign{}
\endlastfoot
S1 CI and PRE-S1 & Passed at exact S1 \texttt{blinded-S1} & Registered tests, source build, and ordinary/strong smoke execution at that identity & Strategy costs, complete-reference coverage, or formal artifact admission \\
Registration and S2 & Descriptive archive reinspected; data-only anchor installed at \texttt{blinded-S2} & Exact S1/S2 identity and closed-Behavior-Set compatibility & Qualification GO or a replayable dispatch authority \\
Qualification & One exact run; q1 completed, q2 was cancelled at the frozen stop, q5 never started & Preregistered Route C decision and fail-closed provenance & Any simulator/native performance estimator or partial formal result \\
Follow-up controls & Five fresh exact-S2 controls passed with authority false & Follow-up source/anchor identity and prerequisite closure & Qualification GO, formal authority, or scientific execution \\
Follow-up qualification & One provider run stopped during hosted-runner setup before checkout or seed admission; zero artifacts & Terminal one-shot disposition and control-plane chronology & Strategy behavior, performance, or a precise watcher root cause \\
Current-source conformance & One exact tagged run and independently rehashed 19-file artifact passed & Eight-case/35-record contract and one fixed whole-query vector & Admission, security, deployment, performance, speedup, or general correctness \\
Primary Route A formal campaign & Not dispatched & The absence of unauthorized execution & Strategy ranking, speedup, ordered-event findings, or native resource claims \\
Validation scaling & One exact tagged attempt; seven jobs/artifacts and the independently rebuilt 54-cell aggregate passed & Current-source semantic equality, compile-call bounds, and descriptive validation-lifecycle scaling & Strategy superiority, native/deployment performance, qualification GO, security, or general correctness \\
\end{longtable}
}

The anonymous supplement records compact stopping provenance, blinded provider
bindings, exact authority states, full numerical tables, and closed claim
mappings for both stopped lineages and the admitted validation-scaling result.
It is part of the circulation packet rather than an alternative source of
empirical results.

\subsection{Synthetic and ordered-event costs}\label{synthetic-and-ordered-event-costs}

The six synthetic and four ordered-event formal shards were not dispatched. Consequently this paper presents no strategy ranking, speedup, projected timing, or cost trade-off plot. The frozen schemas and analysis code remain part of the reproducibility package, but an empty or partial table is not treated as a measurement result.

\subsection{Native OpenFHE execution}\label{native-openfhe-execution}

The six formal native cases were not dispatched. PRE-S1 ordinary and strong smokes demonstrate only that the pinned OpenFHE paths executed at the frozen source identity; they are not benchmark observations. No native median, range, speedup, memory result, or serialized-package comparison is claimed.

\subsection{Preregistered Route C disposition}\label{preregistered-route-c-disposition}

The qualification deadline failure is one of the preregistered falsifiers. It selects Route C without a threshold change, selective rerun, smaller matrix, or reuse of the partial handoff. GitHub's terminal metadata also reported two never-executed downstream cancelled jobs with \texttt{completedAt} one second earlier than \texttt{startedAt}; the controller rejected that terminal observation fail-closed. This provider-boundary anomaly did not create a false GO and did not cause the deadline failure. It is recorded as an implementation limitation for any future lineage, not as permission to modify S1 and repeat this one-shot attempt.

\subsection{Separately preregistered follow-up disposition}\label{separately-preregistered-follow-up-disposition}

The follow-up passed its five prerequisite controls and then consumed its sole
nonserializable capability to dispatch \texttt{blinded-followup-qualification}.
The controller failed closed while establishing the watcher-admission binding
and cancelled the run. The provider recorded only its internal setup step; no
checkout, identity check, seed normalization, registered-seed computation,
producer, replay, guard, native stage, or artifact occurred.

Because the frozen attempt definition attached to the authorized provider
dispatch rather than the first scientific step, that terminal non-success
exhausted the follow-up. Independent reviews rejected both direct rerun and a
repaired watcher on a replacement lineage. The nested provider exception was
not retained, so the exact failing substep remains unknown. This is a
control-plane limitation, not evidence about a maintenance strategy.

\subsection{Current-source deterministic conformance outcome}\label{current-source-deterministic-conformance-outcome}

The independent controller accepted \texttt{blinded-E4-run} at
\texttt{blinded-E4-source}. Every provider-visible step succeeded. The sole
artifact \texttt{blinded-E4-artifact} independently reproduced its provider
digest, contained exactly the 19 preregistered files, passed strict checksum and
JSON-pointer checks, and rehashed every provenance-listed source against a
fresh detached checkout.

At that source, the pinned OpenFHE whole-query witness decrypted to
\texttt{[(0, 128), (4095, 5)]} and matched both independent plaintext oracles;
all 35 deterministic contract records passed. This is a single-source,
single-fixture, version-bound functional observation. It is not candidate
admission, a complete-reference claim, mixed-circuit safety, a security or
deployment result, a performance measurement, or evidence that the stopped
strategy comparison would have succeeded.

\textbf{Table 4. Additional bounded outcomes under anonymous review tokens.}

{\small\def\LTcaptype{none}
\begin{longtable}[]{@{}
  >{\raggedright\arraybackslash}p{(\linewidth - 4\tabcolsep) * \real{0.2500}}
  >{\raggedright\arraybackslash}p{(\linewidth - 4\tabcolsep) * \real{0.3750}}
  >{\raggedright\arraybackslash}p{(\linewidth - 4\tabcolsep) * \real{0.3750}}@{}}
\toprule\noalign{}
Object & Frozen identity & Terminal observation \\
\midrule\noalign{}
\endhead
\bottomrule\noalign{}
\endlastfoot
Follow-up source/freeze & \texttt{blinded-followup-S1/S2}; five fresh controls & Direct-child data-only transition; all authority fields false \\
Follow-up qualification & \texttt{blinded-followup-}\allowbreak\texttt{qualification} & Cancelled during setup; zero scientific steps and zero artifacts \\
Current-source conformance & \texttt{blinded-E4-source/}\allowbreak\texttt{run/}\allowbreak\texttt{artifact} & 35/35 records; one fixed OpenFHE sparse vector matched two oracles \\
\end{longtable}
}

\subsection{Current-source validation-scaling outcome}\label{current-source-validation-scaling-outcome}

The sole tagged attempt, denoted \texttt{blinded-VS-run}, completed all three
producer jobs, all three independent-replay jobs, and the aggregate job. Its
seven provider artifacts had the exact registered names, retention classes,
run/ref/head bindings, and closed contents. A post-terminal audit independently
rehashed the six provider ZIPs, re-inspected every seed package, regenerated
the 54-cell aggregate, and obtained byte equality with
\texttt{blinded-VS-aggregate}.

All 54 compile-call gates passed. The observed count pairs were
\((Q,C)=(5,10),(51,102),(512,1024)\), and all 27 replay cells matched their
bound producer semantic projections. Table~5 reports the registered wall-time
summaries. Each timing entry is median [minimum, maximum] seconds across the
three seeds, rounded half-even to three decimals for space. In the fit column,
\(\alpha\), \(\beta\), and \(R^2\) are rounded half-even to three, four, and
six decimals, respectively. Exact integer
nanoseconds, process times, stage fields, nine-decimal renderings, and reduced
rational fits are in the anonymous supplement. PR-1 denotes
\texttt{periodic-repack/windows=1}, Pad denotes \texttt{padding-reuse}, and
Seg-128 denotes the fixed-width segmented-delta path; P/R denote producer and
independent replay.

\textbf{Table 5. Descriptive current-source validation lifecycle.}

{\scriptsize\setlength{\tabcolsep}{1pt}\def\LTcaptype{none}
\begin{longtable}[]{@{}
  >{\raggedright\arraybackslash}p{0.075\linewidth}
  >{\centering\arraybackslash}p{0.04\linewidth}
  >{\raggedleft\arraybackslash}p{0.17\linewidth}
  >{\raggedleft\arraybackslash}p{0.185\linewidth}
  >{\raggedleft\arraybackslash}p{0.22\linewidth}
  >{\raggedleft\arraybackslash}p{0.24\linewidth}@{}}
\toprule\noalign{}
Path & Role & \(Q=5\) & \(Q=51\) & \(Q=512\) & \(\alpha/\beta/R^2\) \\
\midrule\noalign{}
\endhead
\bottomrule\noalign{}
\endlastfoot
PR-1 & P & 0.507[0.494,0.533] & 3.300[3.248,3.477] & 32.915[32.545,35.868] & 0.114/0.0641/0.999982 \\
PR-1 & R & 0.535[0.463,0.544] & 3.492[2.984,3.515] & 34.046[29.167,34.302] & 0.163/0.0662/0.999994 \\
Pad & P & 0.486[0.474,0.487] & 3.085[3.040,3.214] & 29.918[29.681,31.403] & 0.160/0.0581/0.999995 \\
Pad & R & 0.504[0.435,0.505] & 3.169[2.823,3.247] & 30.843[26.839,30.980] & 0.161/0.0599/0.999993 \\
Seg-128 & P & 1.014[0.976,1.046] & 7.094[6.883,7.358] & 69.763[67.589,71.420] & 0.257/0.1357/0.999995 \\
Seg-128 & R & 1.054[0.976,1.127] & 7.527[6.864,7.929] & 73.518[66.060,77.086] & 0.289/0.1430/0.999998 \\
\end{longtable}
}

The fitted slopes span 0.0581--0.1430 seconds/query and the descriptive
\(R^2\) values span 0.999982--0.999998 on these three points. Those values
describe validation work for this exact source and matrix only. The table is
not a strategy ranking: no pairwise test, speedup estimand, native encrypted
kernel timing, deployment measurement, or population inference was registered.

\section{Limitations}\label{limitations}

The threat model is narrow and provides no formal security theorem. Client B
learns layout and reconstruction metadata; the Cloud learns shapes, counts,
schedules, timing, opaque identifiers, and digests. The historical and
current-source OpenFHE witnesses each cover one frozen fixture and one segment
width. The stopped real-stream matrix uses one source, two deterministic
partitions, and two semantics; its formal strategy-cost and native cases were
never authorized. Neither partial qualification may be reused or rerun.

The admitted scaling study is also narrow: one reviewed source, the S scale,
three fixed query counts, three deterministic seeds, and hosted Ubuntu runners.
It measures validation-bearing Python lifecycle operations, not native
OpenFHE, network service, deployment latency, throughput, or noise safety.
Producer and replay ran in distinct jobs, so their ratio is not a controlled
speedup. Three points make the OLS fits descriptive rather than asymptotic or
inferential; no ranking or population claim follows. Exact provider intervals
are provenance, not fit inputs. The result does not generalize across matrices,
parameters, implementations, workloads, platforms, or adversaries.

\section{Conclusion}\label{conclusion}

Mutable encrypted sparse computation requires more than an updatable container:
matrix state, query reorganization, encrypted execution, and reconstruction
must commit to the same version and be evaluated under complete causal costs.
This work supplies that explicit contract around CSSC, three fixed maintenance
mechanisms, and a fail-closed evidence boundary. The stopped qualification
lineages remain NO-GO for comparative strategy costs. At separately reviewed
current source, however, one fixed OpenFHE witness passed and a one-attempt
54-cell study preserved producer/replay semantics, satisfied every compile-call
bound, and yielded bounded descriptive validation-scaling observations. These
results support version-bound functional and validation-lifecycle claims, not
strategy superiority, native/deployment performance, general correctness, or
formal security.
